\section{Bellman 方程与值迭代、策略迭代}

上一节定义了值函数，却没有告诉我们怎样计算。直接枚举所有未来轨迹会遭遇分支爆炸：每一步有多个动作，每个动作又可能到达多个下一状态，深度增加一层，路径数就成倍增长。动态规划的突破与 HJB 方程相同：不保存整段未来，只把``一步奖励''与``下一状态的最优剩余价值''相加。

\subsection{回报自己带着递归}

折扣回报满足一个纯代数恒等式：

\[
G_t
=R_{t+1}+\gamma G_{t+1}.
\]

在策略 \(\pi\) 下对它取条件期望，并按下一动作与下一状态展开，得到 Bellman 期望方程

\[
V^\pi(s)
=
\sum_a\pi(a\mid s)
\sum_{s'}P(s'\mid s,a)
\left[
r(s,a,s')+\gamma V^\pi(s')
\right].
\]

右端包含 \(V^\pi\) 自己，所以这不是一次向前计算，而是一组联立不动点方程。矩阵形式更清楚。令策略诱导的转移矩阵为 \(P^\pi\)、期望即时奖励向量为 \(r^\pi\)，则

\[
V^\pi=r^\pi+\gamma P^\pi V^\pi,
\qquad
(I-\gamma P^\pi)V^\pi=r^\pi.
\]

有限状态且 \(\gamma<1\) 时，

\[
V^\pi=(I-\gamma P^\pi)^{-1}r^\pi.
\]

这叫策略评估：策略固定后，决策问题退化成 Markov 链上的线性方程。显式求逆不是实际计算的首选，但公式证明了解唯一存在。

\subsection{最优性方程：在第一步取最大}

最优状态值

\[
V^*(s)=\sup_\pi V^\pi(s)
\]

满足 Bellman 最优性方程

\[
V^*(s)
=
\max_a
\sum_{s'}P(s'\mid s,a)
\left[
r(s,a,s')+\gamma V^*(s')
\right].
\]

理由与连续时间 HJB 完全一致：第一步选动作 \(a\)，得到即时奖励并落到 \(s'\)；从 \(s'\) 开始，最优策略的尾巴仍必须最优。若右端由动作 \(a^*(s)\) 取得，取

\[
\pi^*(s)=a^*(s)
\]

就得到一个最优确定性策略。

动作值版本把第一步的动作也留在状态里：

\[
Q^*(s,a)
=
\sum_{s'}P(s'\mid s,a)
\left[
r(s,a,s')+\gamma\max_{a'}Q^*(s',a')
\right].
\]

只要知道 \(Q^*\)，便可直接选 \(\argmax_aQ^*(s,a)\)，不再需要转移模型。这正是下一篇选择学习 \(Q\) 而非只学 \(V\) 的原因。

\subsection{一个两状态例子：长期价值怎样反转眼前偏好}

设状态 \(s\) 有两个动作。动作``取走''立即得 \(2\) 并进入终止状态；动作``等待''立即得 \(0\)，确定到达状态 \(g\)。在 \(g\) 中每一步得 \(1\) 并留在原地。若 \(\gamma=0.9\)，则

\[
V^*(g)=1+0.9V^*(g)=10.
\]

所以在 \(s\) 等待的动作值为 \(0+0.9\times10=9\)，远大于立即取走的 \(2\)。单步奖励偏爱``取走''，长期回报却偏爱``等待''。Bellman 方程没有凭空制造远见，它只是把未来已算出的价值折回当前一步。

若把折扣降到 \(\gamma=0.1\)，则 \(V^*(g)=1/(1-0.1)\)，等待价值约 \(0.11\)，立即取走重新占优。折扣率改变了任务的时间尺度，最优策略随之改变。

\subsection{值迭代：反复应用最优性算子}

定义 Bellman 最优性算子

\[
(TV)(s)
=
\max_a\sum_{s'}P(s'\mid s,a)
\left[r(s,a,s')+\gamma V(s')\right].
\]

值迭代从任意有界初值 \(V_0\) 出发，重复

\[
V_{k+1}=TV_k.
\]

为什么它会收敛，而不是在数值间来回振荡？对任意两个有界函数 \(V,W\)，最大值操作不会放大各动作值之间的最大差异，因此

\[
\norm{TV-TW}_\infty
\le
\gamma\norm{V-W}_\infty.
\]

\(T\) 是 sup 范数下模为 \(\gamma\) 的压缩映射。Banach 不动点定理于是给出：

\begin{itemize}
\tightlist
\item
  \(T\) 有唯一不动点，正是 \(V^*\)；
\item
  从任意有界 \(V_0\) 出发都收敛到 \(V^*\)；
\item
  误差满足 \(\norm{V_k-V^*}_\infty\le\gamma^k\norm{V_0-V^*}_\infty\)。
\end{itemize}

折扣率越接近 \(1\)，压缩越弱，远期影响传播所需迭代越多。值迭代每一轮只把信息向前传一个有效时间步；长规划时域带来的计算成本不会消失。

实际停止时通常检查 Bellman 残差 \(\norm{TV_k-V_k}_\infty\)。残差小能推出到不动点的误差受控，而单看相邻两轮某几个状态变化小不够可靠。数值迭代给近似值后，再对右端取贪心得到策略。

\subsection{策略迭代：先评估，再改进}

值迭代在每轮只做一次 Bellman 备份。策略迭代把工作拆得更明确：

\begin{enumerate}
\def\labelenumi{\arabic{enumi}.}
\tightlist
\item
  策略评估：解 \(V^\pi=r^\pi+\gamma P^\pi V^\pi\)；
\item
  策略改进：对每个状态取
  \[
  \pi_{\mathrm{new}}(s)
  \in
  \argmax_a
  \sum_{s'}P(s'\mid s,a)
  \left[r(s,a,s')+\gamma V^\pi(s')\right].
  \]
\end{enumerate}

策略改进定理保证新策略不会更差。直觉是：旧策略的价值已经把``之后继续照旧''的未来算完；若当前第一步改成一个更高的一步前瞻动作，至少不会降低回报。对有限状态动作集合，若策略不再改变，它就满足 Bellman 最优性方程，因而已经最优；若改变，至少某处严格改善，所以确定性策略有限个意味着过程有限步终止。

完整策略评估可能很贵，可以只做若干轮迭代评估就改进，得到修正策略迭代。值迭代可看成它的极端：每次只做一轮评估便立即改进。两者不是互斥算法，而是一条``每次把当前策略评估多准确''的连续谱。

\subsection{动态规划的前提与边界}

这些算法的保证依赖几项清楚的条件：

\begin{itemize}
\tightlist
\item
  模型已知：每次 Bellman 备份都能对 \(P(s'\mid s,a)\) 精确求期望；
\item
  状态动作可枚举或期望与最大化可计算；
\item
  环境满足 Markov 性且规律稳定；
\item
  折扣情形取 \(\gamma<1\)，或终止任务满足另外的适当条件。
\end{itemize}

状态数巨大时，逐状态扫描遇到离散版的维数灾难；连续状态下更不能直接存一张表。动作连续时，\(\max_a\) 本身也可能是困难的非凸优化。动态规划给出正确方程，却不承诺方程在现实规模上可解。

下一篇只拿走第一项：模型未知，但可以采样。一次样本的

\[
R_{t+1}+\gamma\max_{a'}Q(S_{t+1},a')
\]

会成为 Bellman 右端期望的有噪声观测。把确定性不动点迭代改成随机逼近，就得到 Q 学习。方程没变，能利用的信息变了。
